% 77-07(77쪽) — 07번 그림 · y=f(x) 의 그래프(x=-2, x=2 에서 x축에 뾰족하게 닿고 그 사이는 y축에서 2 인 위로 볼록한 호 · 바깥쪽은 위로 올라감). 스캔 화질이 낮아 TikZ 로 다시 그림.
% 원문에 있는 것만: 눈금 -2, 2(x축), 2(y축) · 라벨 y=f(x), O, x, y (원문에 안내 점선은 없다).
\begin{tikzpicture}[x=0.31cm, y=0.31cm, line width=0.6pt, font=\small]
  \draw[->, >=stealth, line width=0.7pt] (-4.6,0) -- (4.7,0) node[below=1pt] {$x$};
  \draw[->, >=stealth, line width=0.7pt] (0,-1.8) -- (0,5.9) node[left=1pt] {$y$};
  \draw[line width=0.6pt, domain=-3.6:-2, samples=40, smooth] plot (\x, {(\x*\x-4)/2});
  \draw[line width=0.6pt, domain=-2:2, samples=60, smooth] plot (\x, {(4-\x*\x)/2});
  \draw[line width=0.6pt, domain=2:3.6, samples=40, smooth] plot (\x, {(\x*\x-4)/2});
  \node[below, inner sep=2pt] at (-2,0) {$-2$};
  \node[below left=-2pt and -3pt] at (0,0) {$\pt{O}$};
  \node[below, inner sep=2pt] at (2,0) {$2$};
  \node[left, inner sep=2pt] at (0,2) {$2$};
  \node[anchor=west, inner sep=1pt] at (1.3,5.2) {$y=f(x)$};
\end{tikzpicture}
